% GENERATED FILE. Edit canonical lessons, then run npm run generate:books.
% canonical-source-hash: fnv1a64:0dede1d08937ff71
% canonical-lessons: FA-C213-nakhongir, FA-C213-chaharpaye, FA-C213-taqche, FA-C213-dayere, FA-C213-noqte

\chapter{Nail clippers, Stool, Wall niche, Circle, Dot}
\label{ch:fa-nail-clippers-stool-wall-niche-circle-dot}
\hlchaptermodality{213}

\begin{chapteropening}
\textbf{By the end of this chapter:} \emph{I can say nail clippers, a stool, a wall niche, a circle and a dot.}

Complete the last lesson of chapter 213: I can say nail clippers, a stool, a wall niche, a circle and a dot.
\end{chapteropening}

\section[\texorpdfstring{\fa{ناخنگیر} (nâkhongir)}{ناخنگیر (nâkhongir)}]{\fa{ناخنگیر} (nâkhongir) — nail clippers}

\label{lesson:FA-C213-nakhongir}



\begin{quote}
Before the new one: say the Persian for a clay jar, then the Persian for a small tea glass.
\end{quote}

\subsection*{You'll want to know: \fa{ناخنگیر}}
\textbf{\fa{ناخنگیر}} — \emph{nâkhongir} — \textquotedblleft{}nail clippers\textquotedblright{}.

\begin{cousinweb}
\emph{nâkhongir} is \emph{nâkhon}, \textquotedblleft{}nail\textquotedblright{}, with \emph{-gir}, \textquotedblleft{}taking\textquotedblright{} (from \emph{gereftan}, \textquotedblleft{}to take\textquotedblright{}): the nail-taker. Both are native; \emph{nâkhon} is usually traced to Middle Persian \emph{nāxun}.
\end{cousinweb}

\begin{grammarlens}[title={What you've built}]
One more word for this chapter.
\end{grammarlens}

\subsection*{Your turn}
\begin{itemize}
\raggedright
  \item \emph{Say it:} \emph{nâkhongir}
  \item \emph{Say it:} \emph{nâkhongir}, once more
  \item \emph{Say it:} say \emph{nâkhongir}
  \item \emph{Recall:} say the Persian for a clay jar, then the Persian for a small tea glass, then say \emph{nâkhongir} again
\end{itemize}

\begin{tcolorbox}[breakable,colback=teal!4,colframe=teal!35!black,title={Before you move on}]
What does \fa{ناخنگیر} mean? (\textquotedblleft{}Nail clippers\textquotedblright{}.) Say it once more.
\end{tcolorbox}
\section[\texorpdfstring{\fa{چهارپایه} (chahârpâye)}{چهارپایه (chahârpâye)}]{\fa{چهارپایه} (chahârpâye) — a stool}

\label{lesson:FA-C213-chaharpaye}



\begin{quote}
Before the new one: say the Persian for a small tea glass, then the Persian for nail clippers.
\end{quote}

\subsection*{You'll want to know: \fa{چهارپایه}}
\textbf{\fa{چهارپایه}} — \emph{chahârpâye} — \textquotedblleft{}a stool\textquotedblright{}.

\begin{cousinweb}
\emph{chahârpâye} is \emph{chahâr}, \textquotedblleft{}four\textquotedblright{}, with \emph{pâ}, \textquotedblleft{}foot\textquotedblright{}, and \emph{-ye}: a four-footed thing. Both are inherited Persian words.
\end{cousinweb}

\begin{grammarlens}[title={What you've built}]
One more word for this chapter.
\end{grammarlens}

\subsection*{Your turn}
\begin{itemize}
\raggedright
  \item \emph{Say it:} \emph{chahârpâye}
  \item \emph{Say it:} \emph{chahârpâye}, once more
  \item \emph{Say it:} say \emph{chahârpâye}
  \item \emph{Recall:} say the Persian for a small tea glass, then the Persian for nail clippers, then say \emph{chahârpâye} again
\end{itemize}

\begin{tcolorbox}[breakable,colback=teal!4,colframe=teal!35!black,title={Before you move on}]
What does \fa{چهارپایه} mean? (\textquotedblleft{}A stool\textquotedblright{}.) Say it once more.
\end{tcolorbox}
\section[\texorpdfstring{\fa{طاقچه} (tâqche)}{طاقچه (tâqche)}]{\fa{طاقچه} (tâqche) — a wall niche, a shelf}

\label{lesson:FA-C213-taqche}



\begin{quote}
Before the new one: say the Persian for nail clippers, then the Persian for a stool.
\end{quote}

\subsection*{You'll want to know: \fa{طاقچه}}
\textbf{\fa{طاقچه}} — \emph{tâqche} — \textquotedblleft{}a wall niche, a shelf\textquotedblright{}.

\begin{cousinweb}
\emph{tâqche} is \emph{tâq}, \textquotedblleft{}arch, vault\textquotedblright{}, with the small-thing ending \emph{-che}: a little arch. \emph{tâq} is native, usually traced to Middle Persian \emph{tāg}.
\end{cousinweb}

\begin{grammarlens}[title={What you've built}]
One more word for this chapter.
\end{grammarlens}

\subsection*{Your turn}
\begin{itemize}
\raggedright
  \item \emph{Say it:} \emph{tâqche}
  \item \emph{Say it:} \emph{tâqche}, once more
  \item \emph{Say it:} say \emph{tâqche}
  \item \emph{Recall:} say the Persian for nail clippers, then the Persian for a stool, then say \emph{tâqche} again
\end{itemize}

\begin{tcolorbox}[breakable,colback=teal!4,colframe=teal!35!black,title={Before you move on}]
What does \fa{طاقچه} mean? (\textquotedblleft{}A wall niche, a shelf\textquotedblright{}.) Say it once more.
\end{tcolorbox}
\section[\texorpdfstring{\fa{دایره} (dâyere)}{دایره (dâyere)}]{\fa{دایره} (dâyere) — a circle}

\label{lesson:FA-C213-dayere}



\begin{quote}
Before the new one: say the Persian for a stool, then the Persian for a wall niche.
\end{quote}

\subsection*{You'll want to know: \fa{دایره}}
\textbf{\fa{دایره}} — \emph{dâyere} — \textquotedblleft{}a circle\textquotedblright{}.

\begin{cousinweb}
\emph{dâyere} is borrowed from Arabic \emph{dāʾira}, \textquotedblleft{}circle\textquotedblright{}. Its Arabic root is \emph{d-w-r} (Lane).
\end{cousinweb}

\begin{grammarlens}[title={What you've built}]
One more word for this chapter.
\end{grammarlens}

\subsection*{Your turn}
\begin{itemize}
\raggedright
  \item \emph{Say it:} \emph{dâyere}
  \item \emph{Say it:} \emph{dâyere}, once more
  \item \emph{Say it:} say \emph{dâyere}
  \item \emph{Recall:} say the Persian for a stool, then the Persian for a wall niche, then say \emph{dâyere} again
\end{itemize}

\begin{tcolorbox}[breakable,colback=teal!4,colframe=teal!35!black,title={Before you move on}]
What does \fa{دایره} mean? (\textquotedblleft{}A circle\textquotedblright{}.) Say it once more.
\end{tcolorbox}
\section[\texorpdfstring{\fa{نقطه} (noqte)}{نقطه (noqte)}]{\fa{نقطه} (noqte) — a dot, a point}

\label{lesson:FA-C213-noqte}



\begin{quote}
Before the new one: say the Persian for a wall niche, then the Persian for a circle.
\end{quote}

\subsection*{You'll want to know: \fa{نقطه}}
\textbf{\fa{نقطه}} — \emph{noqte} — \textquotedblleft{}a dot, a point\textquotedblright{}.

\begin{cousinweb}
\emph{noqte} is borrowed from Arabic \emph{nuqṭa}, \textquotedblleft{}dot, point\textquotedblright{}. Its Arabic root is \emph{n-q-ṭ} (Lane).
\end{cousinweb}

\begin{grammarlens}[title={What you've built}]
One more word for this chapter.
\end{grammarlens}

\subsection*{Your turn}
\begin{itemize}
\raggedright
  \item \emph{Say it:} \emph{noqte}
  \item \emph{Say it:} \emph{noqte}, once more
  \item \emph{Say it:} say \emph{noqte}
  \item \emph{Recall:} say the Persian for a wall niche, then the Persian for a circle, then say \emph{noqte} again
\end{itemize}

\begin{tcolorbox}[breakable,colback=teal!4,colframe=teal!35!black,title={Before you move on}]
What does \fa{نقطه} mean? (\textquotedblleft{}A dot, a point\textquotedblright{}.) Say it once more.
\end{tcolorbox}
